% Deutsche Parallelfassung zu state/world_to_track.tex

\chapter{World--Track-Koordinatentransformation}
\label{chap:world-to-track}
\index{Projektion!World--Track-Koordinatentransformation}
\index{Bezugssystem!Track Space}

\section*{Motivation}

Dieses Kapitel beantwortet, wie intrinsische Eisenbahnkoordinaten und externe
World-Koordinaten in AIM aufeinander bezogen werden, ohne Geometrie,
Laufzeitbewertung und Koordinatenrepräsentation in einer Operation
zusammenfallen zu lassen. Es begründet, warum die World--Track-Transformation
als Laufzeitdienst auf Operatorebene behandelt wird und was daraus für
Projektion, Mehrdeutigkeit und ingenieurtechnische Interpretation folgt.

Die erste Frage betrifft die Richtung: Was genau wird von World nach Track und
was von Track nach World abgebildet?

\begin{figure}[htbp]
\centering
\begin{tikzpicture}[
  node distance=1.25cm,
  box/.style={draw, rounded corners, minimum width=3.4cm, minimum height=0.8cm, align=center},
  op/.style={draw, rounded corners, minimum width=3.9cm, minimum height=0.85cm, align=center},
  arrow/.style={->, thick}
]
  \node[box] (world) {World Space};
  \node[op, below=of world] (opw) {World-to-Track-Operator};
  \node[box, below=of opw] (track) {Track Space};
  \draw[arrow] (world) -- (opw);
  \draw[arrow] (opw) -- (track);
\end{tikzpicture}

\caption{Unterstützende Operatorsicht: World-to-Track-Transformation als
gerichtete Laufzeitabbildung.}
\label{fig:world-to-track-transformation}
\end{figure}

Die Abbildung legt diese Richtung vor den formalen Definitionen fest und
verhindert die Verwechslung von Projektion mit allgemeiner
Koordinatenkonvertierung.

Eisenbahn-Alignments werden intrinsisch mit bogenlängenbezogenen Koordinaten
formuliert. Reale Daten liegen dagegen üblicherweise in externen
Koordinatenreferenzsystemen wie WGS84, UTM, Gauß--Krüger, DBRef oder lokalen
Ingenieur-CRS vor. Daher ist eine Transformation zwischen intrinsischen
Alignment-Koordinaten und externen World-Koordinaten erforderlich. In AIM
wird sie als Laufzeitbewertungsdienst interpretiert, der intrinsische
Alignment-Zustände, Messdaten, Visualisierung, Realisierung, Topologie und
betriebliche Geometrie verbindet.

\begin{principle}[Primat des intrinsischen Raums]
\label{princ:intrinsic-space-primacy}
\index{Prinzip des intrinsischen Raumprimats}
Innerhalb von AIM ist intrinsische Eisenbahngeometrie primär.
World-Space-Koordinaten werden als externe, durch Laufzeitbewertung erzeugte
Einbettungen interpretiert.
\end{principle}

\section{Gleiskoordinatensystem}

\begin{definition}[Track Space]
\label{def:track-space}
\index{Track Space}
Der Track Space ist die intrinsische Eisenbahnkoordinatendomäne
\[
\Omega_{\mathrm{track}}=\{(s,d,h)\}.
\]
Er repräsentiert Positionen relativ zu einem Alignment statt in einem
externen Koordinatenreferenzsystem.
\end{definition}

\begin{definition}[Gleiskoordinatensystem]
\label{def:track-coordinate-system}
\index{Koordinatensystem!Gleis}
Das intrinsische Gleiskoordinatensystem besteht aus \((s,d,h)\), wobei
\(s\) die intrinsische Bogenlänge, \(d\) den seitlichen und \(h\) den
vertikalen Versatz bezeichnet.
\end{definition}

Gleiskoordinaten gehören zur in \cref{def:track-space} eingeführten
intrinsischen Domäne \(\Omega_{\mathrm{track}}\). In ebenen Situationen
genügt \((s,d)\). Sie dürfen nicht mit projizierten World-Koordinaten
verwechselt werden.

\subsection{Intrinsische Geometrie und externe Koordinaten}

Intrinsische Geometrie ist von jedem externen CRS unabhängig. Dasselbe
intrinsische Alignment kann in projizierte Ingenieurebenen, ellipsoidische
World-Modelle, lokale Realisierungs-Frames oder Visualisierungssysteme
eingebettet werden. Intrinsische Alignment-Geometrie und externe
CRS-Repräsentation bilden damit getrennte Begriffsschichten.

\subsection{Bogenlänge und Stationierung}

\index{Bogenlänge!und Stationierung}
Intrinsische Bogenlänge ist nicht mit ingenieurtechnischer Stationierung
identisch. Betriebliche Kilometersysteme können Kilometersprünge,
Stationsgleichungen, überlappende Bereiche, topologieabhängige Bezüge und
betriebliche Zweigbeziehungen enthalten. Stationierung ist eine betriebliche
Referenzstruktur über der intrinsischen Geometrie; die Abbildung
\(s\leftrightarrow\mathrm{km}\) ist im Allgemeinen weder linear noch global
eindeutig.

\section{Track-to-World-Transformation}

\begin{definition}[Track-to-World-Abbildung]
\label{def:state-track-to-world-runtime}
\index{Projektion!Track-to-World}
Die Vorwärtstransformation lautet
\begin{equation}
x(s,d)=\gamma(s)+d\,\mathbf n(s).
\label{eq:state-track-to-world-runtime}
\end{equation}
\end{definition}

Diese Beziehung spezialisiert den in
\cref{def:track-to-world-operator} eingeführten kanonischen Operator
\(\mathcal P_{\mathrm{tw}}\).

Die ergänzende Frage lautet: Wie wird der intrinsische Zustand in den World
Space ausgegeben, nachdem die Richtung feststeht?

\begin{figure}[htbp]
\centering
\begin{tikzpicture}[
  node distance=1.25cm,
  box/.style={draw, rounded corners, minimum width=3.4cm, minimum height=0.8cm, align=center},
  op/.style={draw, rounded corners, minimum width=3.9cm, minimum height=0.85cm, align=center},
  arrow/.style={->, thick}
]
  \node[box] (track) {Track Space};
  \node[op, below=of track] (opt) {Track-to-World-Operator};
  \node[box, below=of opt] (world) {World Space};
  \draw[arrow] (track) -- (opt);
  \draw[arrow] (opt) -- (world);
\end{tikzpicture}

\caption{Unterstützende Operatorsicht: Track-to-World-Transformation als
ergänzende gerichtete Abbildung.}
\label{fig:track-to-world-transformation-support}
\end{figure}

Diese Abbildung schließt zusammen mit
\cref{fig:world-to-track-transformation} das Operatorpaar. Beide Abbildungen
müssen als gekoppelte Laufzeitdienste mit unterschiedlichem numerischem
Verhalten behandelt werden, nicht als ein Koordinatenwerkzeug.

\subsection{Ebene Tangenten- und Normalenrichtungen}

Im ebenen Fall gilt
\[
\mathbf t(s)=(\cos\theta(s),\sin\theta(s)),
\qquad
\mathbf n(s)=(-\sin\theta(s),\cos\theta(s)).
\]

\subsection{Räumliche Laufzeitbewertung}

Bei räumlicher Laufzeitbewertung hängt die Abbildung zusätzlich von Profil,
Überhöhung, Frame-Konstruktion, metrischer Interpretation und
Laufzeit-Realisierungskontext ab. Die World-Einbettung \(x(s,d,h)\) wird daher
im Allgemeinen dynamisch aus der Laufzeit-pose3-Bewertung erzeugt.

\subsection{Interpretation der Referenzlinie}

Die Seitenkoordinate \(d\) hängt vom gewählten Bezug ab. Mögliche Strukturen
sind Alignment-Mittellinie, Gleismitte, Schienenachse, Spurweitenreferenz oder
betriebliche Versatzdefinition. Eine formalere Behandlung für Mehrgleissysteme,
überhöhungsabhängige Versätze und Vehicle Space bleibt offen.

\section{World-to-Track-Transformation}

\begin{definition}[World-to-Track-Projektion]
\label{def:state-world-to-track-runtime}
\index{Projektion!World-to-Track}
Die inverse Transformation lautet
\begin{equation}
(s^\ast,d^\ast)
=
\arg\min_{s,d}
\|x-(\gamma(s)+d\,\mathbf n(s))\|.
\label{eq:state-world-to-track-runtime}
\end{equation}
\end{definition}

Diese Beziehung spezialisiert den in
\cref{def:world-to-track-operator} eingeführten Operator
\(\mathcal P_{\mathrm{wt}}\). Der Begriff Projektion wird betrieblich
verwendet und impliziert nicht notwendig eine eindeutige orthogonale Projektion
im klassischen geometrischen Sinn. In AIM ist sie ein inverses
Laufzeitbewertungsproblem.

\section{Koordinatentransformation und Projektion}

\index{Projektion!Unterschied zur Koordinatentransformation}
Koordinatentransformation und World-to-Track-Projektion sind verschiedene
Operationen. Koordinatentransformationen bilden zwischen
World-Space-Repräsentationen ab,
\(\mathcal C:\Omega_{\mathrm{world}}\to\Omega_{\mathrm{world}}\), wie in
\cref{def:coordinate-transform-operator} eingeführt.
World-to-Track projiziert zwischen \(\Omega_{\mathrm{world}}\) und
\(\Omega_{\mathrm{track}}\).

\begin{principle}[Projektionsunterscheidungsprinzip]
\label{princ:projection-distinction}
\index{Projektionsunterscheidungsprinzip}
Koordinatentransformation darf nicht mit intrinsischer
World-to-Track-Projektion verwechselt werden.
\end{principle}

CRS-Transformationen wirken vollständig innerhalb externer
Koordinatenrepräsentationen; Projektion greift auf intrinsische
Eisenbahngeometrie zu.

\section{Numerische Bewertung}

Die inverse Transformation ist im Allgemeinen nichtlinear. Typische Verfahren
sind Nächster-Punkt-Suche, lokale Newton-Iteration, abschnittsweise Projektion,
topologiebewusster Lookup und hybride Laufzeitstrategien. Die Qualität hängt
stark von Krümmungsverhalten, Topologie, Initialisierung,
Laufzeitqualifikationen und metrischer Interpretation ab.

\section{Mehrdeutigkeit und Kontextabhängigkeit}

\index{Projektion!Mehrdeutigkeit}
Die World-to-Track-Abbildung ist im Allgemeinen nicht global eindeutig.
Mehrdeutigkeiten können in Schleifen, Weichen und Kreuzungen, parallelen
Gleisen, eng benachbarten Alignments, Bahnhofsbereichen und topologiereichen
Netzen auftreten. Mehrere intrinsische Zustände können nahezu derselben
World-Position entsprechen.

Projektion ist daher grundlegend kontextabhängig. Sie kann Track- und
Route-Identity, Topologie, erwarteten Stationsbereich, Betriebsrichtung,
Interaktionsgeschichte, Fahrzeugkontext oder Netzzustandsbedingungen
benötigen.

\begin{principle}[Projektionsmehrdeutigkeitsprinzip]
\label{princ:projection-ambiguity}
\index{Projektionsmehrdeutigkeitsprinzip}
World-to-Track-Projektion kann im Allgemeinen nicht als rein lokale
geometrische Operation interpretiert werden.
\end{principle}

\section{Topologie und betriebliche Referenzen}

Intrinsische Geometrie allein reicht für betriebliche Eisenbahnreferenzierung
nicht aus. Die Interpretation hängt zusätzlich von Topologie,
Routenzugehörigkeit, Kilometerbezügen, Betriebsrichtung und Netzsemantik ab.
Ein vollständiger Projektionskontext kann daher
\[
(x,\mathcal T,\mathcal R,\mathcal O)
\]
erfordern, wobei \(x\) die World-Eingabe, \(\mathcal T\) die Topologie,
\(\mathcal R\) Route/Referenz und \(\mathcal O\) den Betriebskontext
bezeichnet.

\section{Beziehung zur Krümmung}

Projektionsstabilität hängt unmittelbar vom Krümmungsverhalten ab. Hohe
Krümmung erhöht die Empfindlichkeit, niedrige Krümmung ergibt üblicherweise
stabiles Verhalten. Fehler in intrinsischen Koordinaten pflanzen sich
nichtlinear in den World Space fort; das Verhalten nahe Übergängen kann stark
von Krümmung höherer Ordnung abhängen.

\section{Laufzeitabhängigkeit}

In AIM hängt die World--Track-Transformation von der Laufzeitbewertung der
zugrunde liegenden Alignment-Repräsentation ab: cGeom definiert intrinsische
Geometrie, pose2 Tangentenentwicklung, pose3 räumliche Interpretation,
metrische Dienste die realisierte Einbettung, CRS-Dienste die externe
Repräsentation und Laufzeitoperatoren die resultierenden Zustände.

Unterschiedliche Konfigurationen können verschiedene Ergebnisse liefern, etwa
projizierte oder ellipsoidische Bewertung, Ziel- oder realisierte Geometrie,
niedrige oder hohe Auflösung sowie vereinfachte oder fahrzeugbewusste
Interpretation.

\begin{principle}[Laufzeitprojektionsprinzip]
\label{princ:runtime-projection}
World--Track-Transformation ist ein auf Alignment-Zuständen arbeitendes
Laufzeitbewertungsproblem und kein statisches Koordinatenwerkzeug.
\end{principle}

\section{LOD und hybride Bewertung}

Exakte World-to-Track-Projektion ist rechenintensiv. Große Systeme benötigen
daher LOD-abhängige und näherungsweise Bewertung, topologiebewusste Filterung
und hybride Laufzeitarchitekturen.

\begin{proposition}[Selektives Projektionsprinzip]
\label{prop:selective-projection}
Exakte Projektion soll nur dort bewertet werden, wo hohe Genauigkeit
erforderlich ist.
\end{proposition}

Eine hybride Strategie verbindet vorberechnete Stichproben für groben Lookup,
topologiebewusste Kandidatenfilterung, zwischengespeicherte Laufzeitzustände
und lokale parametrische Verfeinerung. Dies ist besonders für große Netze,
interaktive Visualisierung, Echtzeit-Vehicle-Space-Bewertung und
mehrskalige Laufzeitsysteme wichtig.

\section{Verbindung zu CRS und Realisierung}

Die vollständige externe Transformationskette lautet
\[
\text{CRS}\rightarrow\Omega_{\mathrm{world}}\rightarrow\Omega_{\mathrm{track}}.
\]
Für hochgenaue Anwendungen kann das Alignment auf einem Referenzellipsoid
statt in einer projizierten Ebene definiert sein. Projektion kann auch von
realisierter Geometrie abhängen:
\[
\gamma_{\mathrm{real}}(s)\neq\gamma_{\mathrm{target}}(s).
\]
Physische Messungen sollen daher auf realisierte statt rein analytische
Geometrie bezogen werden. Die Trennung zwischen intrinsischer Geometrie und
externer Einbettung ist grundlegend.

\section{Zentrale Erkenntnis}

\begin{proposition}
World--Track-Transformation ist ein inverser Laufzeitdienst und keine
statische Koordinatenkonvertierung.
\end{proposition}
\begin{proposition}
In AIM ist der intrinsische Alignment-Raum primär; World-Koordinaten sind
externe Einbettungen von Alignment-Zuständen.
\end{proposition}
\begin{proposition}
Projektion ist im Allgemeinen nichtlinear, mehrdeutig, topologie-,
laufzeit- und kontextabhängig.
\end{proposition}

\section{Verbindung zu AIM}

\begin{summary}
World--Track-Transformation verbindet intrinsische Eisenbahngeometrie mit
externen Koordinatensystemen, Topologie, Laufzeitbewertung und gemessener
Realität. Ihre nichtlineare, mehrdeutige, topologie- und skalenabhängige Natur
macht sie zu einer zentralen Laufzeitkomponente von AIM.

Projektion ist keine statische Koordinatenkonvertierung, sondern kontextuelle
Laufzeitrekonstruktion intrinsischer Eisenbahnzustände aus externen
Beobachtungen. Effiziente Behandlung verlangt hybride Laufzeitarchitekturen,
topologiebewusste Strategien, LOD-abhängige Bewertung, dünn besetzte
Zustandsrekonstruktion und metrikbewusste Laufzeitdienste.
\end{summary}
